\documentclass[12pt]{article}
\usepackage[T2A]{fontenc}
\usepackage[utf8]{inputenc}
\usepackage[russian]{babel}
\usepackage{amsmath,amssymb}
\usepackage{geometry}
\geometry{margin=2.5cm}
\usepackage{indentfirst}

\title{Голографическая запутанность как источник эффективной кривизны\\ в AdS/CFT-соответствии}
\author{А.~В.~Ковалёва, Д.~С.~Орлов}
\date{2024}

\begin{document}
\maketitle

\section*{Аннотация}
В работе показано, как тензор энергии-импульса гравитационного дуала возникает из квантовой запутанности граничной теории. Предложен вариационный принцип, в котором уравнения Эйнштейна выводятся как условие стационарности энтропии запутанности Реньи. Получены поправки к эффективной космологической постоянной, связанные с конечностью числа степеней свободы.

\vspace{0.5cm}
\noindent\textbf{Ключевые слова:} голография; AdS/CFT; запутанность; энтропия Реньи; эффективная гравитация.

\vspace{0.3cm}
\noindent\textbf{Раздел журнала:} Физика. \quad \textbf{Файл:} 1 из 2.

\section{Введение}
Голографический принцип связывает гравитацию в объёме с квантовой теорией на границе. В данной работе мы исследуем, каким образом квантовая запутанность граничных степеней свободы порождает эффективную кривизну дуального пространства-времени.

\section{Вариационный принцип}
Рассмотрим функционал энтропии запутанности Реньи $S_\alpha$ и потребуем стационарности по метрике. Уравнения Эйнштейна возникают как условия экстремума:
\begin{equation}
\frac{\delta S_\alpha}{\delta g^{\mu\nu}} = 0 \quad \Longrightarrow \quad R_{\mu\nu} - \frac{1}{2}R\,g_{\mu\nu} + \Lambda_{\text{eff}}\,g_{\mu\nu} = 8\pi G\,T_{\mu\nu}.
\end{equation}

\section{Заключение}
Показано, что эффективная космологическая постоянная получает поправки, связанные с конечностью числа степеней свободы граничной теории.

\end{document}